% Generated shared reviewed material.
% Actual signed divisor variable; no new asymptotic saving or Lean declaration.
\section{The signed divisor mean and a complete-column laboratory}\label{sd:core}
The preceding signed-slice estimate loses the sign of the common ideal
 divisor. We now retain that variable in an exact analytic reduction and
in a complete finite computation. Neither supplies the missing
$5/12$ saving at $H=D^{3/4}$. The computations below use a new fixed
profile; they do not fill in the uncomputed columns of the earlier
broad-plateau profile.

Let $\mathcal C$ be original squarefree good ideal columns in $[D,2D]$,
and let $\mathcal R$ be an actual finite literal row block, with
nonnegative multiplicities $w_u$. All shared-prime zeros of the original
$\chi_m(u)$ remain. For arbitrary complex coefficient vectors define the
linear-first form
\begin{equation}\label{sd:form}
 DQ_1(x,y)=\sum_{u\in\mathcal R}w_u
 \sum_{\substack{m,n\in\mathcal C\\(m,n)=1}}
 x_m\overline{y_n}\chi_m(u)\overline{\chi_n(u)}.
\end{equation}
The original component is obtained with
$x_m=X_U(m)\nu(m)W(Nm/D)$ and the corresponding $y_m$.
On the nonprincipal row blocks used below the original residue term
vanishes. Equation~\eqref{sd:form} is a raw coprime form on a generic
block; principal rows would still require their separate residue assembly.

For every ideal divisor of an active column put
\[
 F_d(x;u)=\sum_{m\in\mathcal C,\ d\mid m}x_m\chi_m(u),\qquad
 Z_d(x,y)=\sum_u w_uF_d(x;u)\overline{F_d(y;u)}.
\]
The exact identity is
\begin{equation}\label{sd:incidence}
 DQ_1(x,y)=\sum_d\mu(d)Z_d(x,y).
\end{equation}
There is no extra $\mu(q)$: the vectors already contain the original
coefficients. Full divisors $d=m$ are essential to cancel self-pairs.
Divisors and squarefreeness refer to ideals, including distinct ideals
with the same rational norm.

\begin{proposition}[A weighted divisor-incidence lift]\label{sd:lift}
Suppose $\lambda_d\ge0$ is an arbitrary-coefficient bound
\[
 \sum_u w_u|F_d(x;u)|^2\le\lambda_d
                 \sum_{m\in\mathcal C,\ d\mid m}|x_m|^2.
\]
With $B(m)=\sum_{d\mid m}\lambda_d$ and
$\|x\|_B^2=\sum_mB(m)|x_m|^2$, one has
\begin{equation}\label{sd:weighted}
 |DQ_1(x,y)|\le\|x\|_B\|y\|_B.
\end{equation}
For a finite literal phase table one may take the exactly computed
\[
 \lambda_d=E_d:=\sum_u w_u\#\{m\in\mathcal C:d\mid m,\ \chi_m(u)\ne0\}.
\]
This last choice is a Frobenius certificate, not a new analytic estimate.
For $x=t+e$ it gives the finite perturbation bound
\begin{equation}\label{sd:certificate}
 |Q_1(x,x)-Q_1(t,t)|\le
 D^{-1}\bigl(2\|t\|_B\|e\|_B+\|e\|_B^2\bigr).
\end{equation}
\end{proposition}
\begin{proof}
Insert ideal-divisor inversion in~\eqref{sd:form}. Apply row Cauchy to
each $Z_d$ and weighted Cauchy to the resulting sum over $d$.
Then $\sum_d\lambda_d\|x_d\|_2^2=\|x\|_B^2$.
The squared Frobenius norm of the actual $F_d$ matrix is exactly $E_d$.
Expand $Q_1(t+e,t+e)-Q_1(t,t)$ to obtain the last assertion.
Positivity of the coprime-truncated matrix was never assumed.
\end{proof}
Using an accepted positive-row sieve bound for $\lambda_d$ only recovers
the old exponent up to ideal-divisor factors. It does not estimate the
signed sum before the triangle in~\eqref{sd:incidence}.

There is also a genuinely signed scalar expansion. Fix squarefree residual
ideals $a,b$, which may overlap, and an actual row $u$.
Keep the full deleted set $\mathcal T=S\cup\operatorname{supp}(abu)$ and set
\[
 \mathcal Z_{U;a,b,\mathcal T}(s)=
 \sum_{\substack{d\ \mathrm{squarefree}\\(d,\mathcal T)=1}}
 \frac{\mu(d)A_U(da)A_U(db)}{(Nd)^s},\qquad\Re s>1.
\]
For four positive cutoff variables put
$P_p=1-(Np)^{-v_1}-(Np)^{-v_2}$ and
$Q_p=1-(Np)^{-w_1}-(Np)^{-w_2}$.
The common-divisor Euler factor is $1-(Np)^{-s}P_pQ_p$.
Its first-order factorization has four numerator and five denominator
zeta factors:
\begin{equation}\label{sd:nine}
 \prod_{p\notin\mathcal T}(1-(Np)^{-s}P_pQ_p)
 =\frac{\prod_{i=1}^2\zeta_K^{\mathcal T}(s+v_i)
                         \zeta_K^{\mathcal T}(s+w_i)}
 {\zeta_K^{\mathcal T}(s)
       \prod_{i,j=1}^2\zeta_K^{\mathcal T}(s+v_i+w_j)}\,
       \mathcal H_{\mathcal T}(s;v,w).
\end{equation}
The correction is normally convergent at
$\Re s\ge\tfrac12+\eta$, $\Re v_i,\Re w_j\ge0$, for each fixed
$\eta>0$. This says nothing about zeros or residues of the zeta ratio.
Finite-height Perron kernels, with their explicit endpoint/error costs,
recover the cutoffs by successive limits.

A direct finite cutoff expansion gives more explicit information:
\begin{equation}\label{sd:finite-numerator}
 \mathcal Z_{U;a,b,\mathcal T}(s)
       =\frac{\mathcal N_{U;a,b,\mathcal T}(s)}
                    {\zeta_K^{\mathcal T}(s)}.
\end{equation}
The numerator is a finite sum over eight nonzero joint short-divisor
states, holomorphic on $\Re s>0$. Uniformly for $\sigma\ge\sigma_0>0$,
\[
 |\mathcal N_{U;a,b,\mathcal T}(\sigma+i\tau)|
 \ll_{K,\sigma_0,\epsilon}(Na\,Nb)^\epsilon
             U^{4\max(1-\sigma,0)+\epsilon}.
\]
The costs near $\sigma=1/2$ persist. The remaining long state has the
M\"obius sign, but the eight finite short states cannot be omitted.
A fixed-residual reciprocal zeta factor also does not prove an
uncancelled pole in the complete coherent sum.

\paragraph{Exact complete-column reference.}
Choose once $\delta=1/240000$ and
$W_\delta(y)=s((y-1)/\delta)s((2-y)/\delta)$, where $s$ is the smooth
step from zero to one. For the finite grid $D\le60000$ every strict
integer norm node has $W_\delta=1$; endpoints vanish. Its derivatives
cost $O_k(\delta^{-k})$. The three cases have different declared row
blocks, not a common asymptotic scaling experiment. They use principal
$\nu$, $S$ above $2,3$, literal class one and nonprincipal actual rows.
\begin{center}
\begin{tabular}{r|r|r|r|r}
 $D$ & columns & divisors & rows & $S_1^{AA}$\\\hline
 $144$ & $41$ & $53$ & $5$ & $-19/36$\\
 $576$ & $164$ & $212$ & $11$ & $-221/288$\\
 $60000$ & $16953$ & $21750$ & $20$ & $-5119/5000$
\end{tabular}
\end{center}
In the last case the exact even/odd M\"obius masses are
$543416$ and $604844$, giving numerator $-61428$ before division by $D$.
All $339060$ raw column--row values, including $47908$ zeros, and all nine
cross components are retained. In particular $1361$ columns have
nonsquarefree rational norm while remaining squarefree as ideals.
The \href{https://amelie-iska.github.io/GRH-research/signed-divisor.html}{signed-divisor inspector}
exposes the actual states and dimensions. These finite negative slices
are neither whole-family asymptotic counterexamples nor a positive
arithmetic determinant or GRH result.
